% LOMBEIDA ESCALERAS — sustituir marcadores antes de la entrega
\renewcommand{\fuentespie}{Fuentes: \cite{assiri2025,mushtaq2024}}
\begin{frame}{Tolerancia a fallos y tipos de fallos}
  \begin{columns}[T,onlytextwidth]
    \begin{column}{0.46\textwidth}
      \begin{block}{Tolerancia a fallos}
        Seguir funcionando correctamente aunque ocurran fallos.~\cite{assiri2025}
      \end{block}

      \vspace{0.8em}
      {\small
      \textbf{Fallo} (causa)\\
      $\downarrow$ \textbf{error} (estado)\\
      $\downarrow$ \textbf{falla} (servicio)~\cite{mushtaq2024}\par}
    \end{column}

    \begin{column}{0.50\textwidth}
      \begin{block}{Tipos de fallos}
        \small
        \begin{itemize}
          \item \textbf{Nodos:} máquina caída o lenta.
          \item \textbf{Comunicación:} mensajes perdidos o tardíos.
          \item \textbf{Procesos:} bloqueo o respuesta errónea.
          \item \textbf{Recursos:} CPU, memoria o disco agotados.~\cite{mushtaq2024}
        \end{itemize}
      \end{block}
    \end{column}
  \end{columns}
\end{frame}
\renewcommand{\fuentespie}{Fuentes: [PENDIENTE]}

\renewcommand{\fuentespie}{Fuentes: \cite{mushtaq2024,asch2025,chaurasia2024}}
\begin{frame}{Detección y recuperación ante fallos}
  \begin{columns}[T,onlytextwidth]
    \begin{column}{0.46\textwidth}
      \begin{block}{Detección}
        \small
        \begin{itemize}
          \item \textbf{Heartbeats:} señales periódicas de actividad.
          \item \textbf{Timeouts:} límite de espera para sospechar un fallo.~\cite{chaurasia2024}
          \item \textbf{Watchdogs:} vigilan señales tardías o ausentes.~\cite{asch2025}
        \end{itemize}
      \end{block}
    \end{column}

    \begin{column}{0.50\textwidth}
      \begin{block}{Recuperación}
        \small
        \begin{itemize}
          \item \textbf{Replicación:} mantener copias de respaldo.
          \item \textbf{Reasignación:} trasladar tareas a otro nodo.
          \item \textbf{Checkpoint y reinicio:} retomar un estado guardado.~\cite{mushtaq2024}
        \end{itemize}
      \end{block}
    \end{column}
  \end{columns}
\end{frame}
\renewcommand{\fuentespie}{Fuentes: [PENDIENTE]}

\renewcommand{\fuentespie}{Fuentes: \cite{assiri2025,mushtaq2024}; práctica propia}
\begin{frame}[t]{SGRA: lecciones para la PFC y conclusiones}
  \begin{columns}[T,onlytextwidth]
    \begin{column}{0.46\textwidth}
      \begin{block}{Aplicación a SGRA}
        \small
        \begin{itemize}
          \item \textbf{Estudiante:} solicita refuerzo.
          \item \textbf{Java:} comprueba cuatro franjas.
          \item \textbf{Docente:} programa una franja libre.
        \end{itemize}
      \end{block}
    \end{column}

    \begin{column}{0.50\textwidth}
      \begin{block}{Resultados de la práctica}
        \small
        \begin{itemize}
          \item \textbf{Paralelismo:} uno frente a cuatro trabajadores.
          \item \textbf{Fallo simulado:} reintento de una franja.
          \item \textbf{Fallo persistente:} franja sin comprobar.
        \end{itemize}
      \end{block}
    \end{column}
  \end{columns}
  \vspace{1em}
  \begin{beamercolorbox}[wd=\textwidth,sep=0.8em]{block body}
    \small\textbf{Conclusión:} recuperar solo la tarea fallida evita repetir el trabajo ya completado.~\cite{assiri2025,mushtaq2024}
  \end{beamercolorbox}
\end{frame}
\renewcommand{\fuentespie}{Fuentes: [PENDIENTE]}
